\section{The first sharp staircase beyond the residue characteristic}
\label{sec:sharp-range}

The finite algorithm can be complemented by an explicit answer in a range
where the number of terms exceeds the residue characteristic.  Write
\[
 \mathcal D_p=\{x\in\overline{\mathbf Q}_p:
                  v_p(x-1)>1/(p-1)\}.
\]
Coefficients and roots may lie in arbitrary finite extensions of
$\mathbf Q_p$, and zeros are counted with their algebraic multiplicities.
The maximum in the following theorem is over all supports of the indicated
size, as well as over coefficients.

\begin{theorem}[A sharp staircase]\label{thm:sharp-staircase}
Let $p$ be an odd prime and let $p+2\le t\le 2p-1$.  The largest possible
number of zeros in $\mathcal D_p$ of a nonzero polynomial with at most $t$
terms is
\begin{equation}\label{eq:sharp-staircase}
 B_p(t)=p^2+\left\lfloor\frac{t-p-2}{2}\right\rfloor.
\end{equation}
The same number is attained by an integer polynomial with exactly $t$
terms which is irreducible over $\mathbf Q$.  In particular, all its zeros
are simple.  Replacing $x-1$ by $x/A-1$, for any nonzero algebraic
$p$-adic number $A$, leaves the maximum unchanged.
\end{theorem}

For example, the maxima at $p=7$ and $t=9,10,11,12,13$ are respectively
$49,49,50,50,51$.  The first value extends the repository's $p+2$-term
theorem; the successive increases require two additional terms each.
For $p=3$ the range consists only of $t=5$, so the assertion there is the
already-known value nine.  The new staircase begins at $p\ge5$.

We prove a more precise assertion about supports.  If $E$ is a finite set
of distinct nonnegative exponents, put
$s_a=\#\{n\in E:n\equiv a\pmod p\}$.

\begin{theorem}[One large residue class]\label{thm:large-residue-class}
Suppose $|E|\le2p-1$, at least two of the $s_a$ are positive, and one
residue class has multiplicity $m>p$.  Put
$q=\max\{s_a:a\text{ is not that class}\}$ and $d=\min(m-p,q)$.
The exact maximum number of zeros in $\mathcal D_p$ for nonzero
polynomials supported on $E$ is
\begin{equation}\label{eq:large-profile}
                   p^2+d-1.
\end{equation}
Every exponent of $E$ can be required to occur, and an attaining
polynomial can be chosen with rational coefficients and simple nonzero
zeros.  This formula makes no assertion about a support confined to a
single residue class.
\end{theorem}

\subsection{A normalized moment test}

We record the analytic reduction needed in both directions.  For
$f(X)=\sum_{n\in E}c_nX^n$, let $s=|E|$ and
\[
 P(T)=\prod_{n\in E}(T-n),\qquad
 A=\mathbf Z_{(p)}[T]/P(T),\qquad
 M_k=\sum_{n\in E}c_n n^k.
\]
The functional $L(T^k)=M_k$ is well-defined on $A$.  Because $P$ is
monic, every $M_k$ is an integral linear combination of
$M_0,\ldots,M_{s-1}$.  The Vandermonde matrix is invertible over
$\mathbf Q$, so these first $s$ moments cannot all vanish when $f\ne0$.
Divide them by one of least valuation.  All moments are now integral,
and at least one of the first $s$ is a unit.

For rational $r>1/(p-1)$, choose an algebraic $z$ with $v_p(z)=r$.
The substitution $X=\exp(zY)$ identifies the closed unit disk in $Y$
with $v_p(X-1)\ge r$, and preserves multiplicity.  Its coefficients are
$M_kz^k/k!$.  If $s_p(k)$ denotes the base-$p$ digit sum and
$\delta=r-1/(p-1)>0$, their valuations are
\begin{equation}\label{eq:sharp-weight}
 v_p(M_k)+w_r(k),\qquad
 w_r(k)=\delta k+\frac{s_p(k)}{p-1}.
\end{equation}
The number of geometric zeros on the closed disk is the largest index
at which these valuations take their minimum.  This is the
Weierstrass degree of the restricted power series.

Suppose $s\le2p-1$ and such a minimizing index $k$ satisfies $k\ge s$.
Every $i<s$ has $s_p(i)\le p-1$.  Comparison with an early unit moment
gives
\[
 0\le v_p(M_k)
 \le\delta(i-k)+\frac{s_p(i)-s_p(k)}{p-1}.
\]
Consequently
\begin{equation}\label{eq:sharp-small-moment}
               s_p(k)\le p-2,\qquad v_p(M_k)<1.
\end{equation}
In particular, congruences modulo $p$ suffice to exclude a large
minimizing index.  If in $A/pA$ a power $T^k$ is a linear combination
of powers $T^\ell$ with
$w_r(\ell)<w_r(k)$, lift the relation to $A$.  The error has moment
valuation at least one.  By \eqref{eq:sharp-small-moment}, at least one
participating $M_\ell$ has valuation no larger than $v_p(M_k)$; its
series coefficient then has strictly smaller valuation.  This is a
contradiction.

\subsection{The upper bound}

First suppose every residue multiplicity is at most $p$.  In
\[
 R=A/pA\simeq\prod_{a:s_a>0}\mathbf F_p[T]/(T-a)^{s_a}
\]
we have $T^{p^j}=T^p$ for every $j\ge1$.  A power of digit sum at most
$p-2$ therefore has the same image as $T^{i+pj}$ with
$i+j\le p-2$.  If its original index is at least $p^2$, this replacement
strictly decreases the index and does not increase the digit sum.
The normalized moment test excludes it.  Hence in this case there are
fewer than $p^2$ zeros in every closed disk under consideration.

Next suppose one multiplicity $m$ exceeds $p$ and another is positive.
It is unique, and $m\le2p-2$.  Multiplication of $f$ by a suitable
nonnegative power of $X$ allows its residue class to be taken as zero;
this does not change any zero in the disk.  In $R$ put
\[
                  U=T^p,\qquad Z=T^{p^2}.
\]
On every factor other than the zero factor, $U=Z=a$; on the zero factor,
$U=T^p$, $Z=0$, and $U^2=0$.  Every higher Frobenius power equals $Z$.
Thus a candidate of digit sum $h\le p-2$ has image
$T^iU^jZ^\ell$, where $i+j+\ell=h$.

If $\ell=0$, its compressed index is below $p^2$.  If
$\ell\ge1$ and $j+\ell\ge2$, its image equals
$T^iU^{j+\ell}$: this follows factor by factor, using $U^2=Z=0$
on the zero factor.  Again the replacement index is below $p^2$ and
has the same digit sum.  The only remaining candidate above $p^2-1$
is therefore
\begin{equation}\label{eq:sharp-exceptional-index}
                        k=p^2+h-1.
\end{equation}

If $h>m-p$, then $T^{h-1}Z=T^{p+h-1}$ in $R$, since both vanish on
the zero factor.  This is a lower-index replacement.  If $h>q$, then
$Z(T-Z)^{h-1}=0$ in $R$: on every nonzero factor the nilpotent
$T-a$ has order at most $q\le h-1$.  Expanding gives
\[
 T^{h-1}Z
   =\sum_{j=2}^{h}(-1)^j\binom{h-1}{j-1}T^{h-j}Z^j
   =\sum_{j=2}^{h}(-1)^j\binom{h-1}{j-1}T^{h-j}U^j.
\]
Every term on the right has digit sum $h$ and index below $p^2$.
The moment test excludes \eqref{eq:sharp-exceptional-index} unless
$h\le\min(m-p,q)$.  This proves the upper bound in
Theorem~\ref{thm:large-residue-class}.

For the unrestricted theorem a support confined to one residue class
causes no difficulty.  Factor out its smallest monomial and divide all
exponent differences by their largest common power $p^a$.  The
polynomial becomes $X^b g(X^{p^a})$, where $g$ occupies at least two
residue classes.  On $\mathcal D_p$ the power map is injective, because
the logarithm identifies it with multiplication by $p^a$, and it
preserves multiplicities.  Its image is contained in $\mathcal D_p$.
Thus the bound for $g$ is also a bound for the original polynomial.
The case of a single monomial has no zeros in the disk.

For any remaining support of size at most $t$, either the bound is
less than $p^2$, or
$2\min(m-p,q)\le(m-p)+q\le t-p$.  Therefore
\[
 p^2+\min(m-p,q)-1
 \le p^2+\left\lfloor\frac{t-p}{2}\right\rfloor-1,
\]
which is \eqref{eq:sharp-staircase}.  Every polynomial has finitely many
roots, so bounds on every rational closed disk of depth greater than
$1/(p-1)$ also bound the roots in the open disk $\mathcal D_p$.

\subsection{A functional which attains the bound}

Let $1\le d\le p-2$ and suppose that the support has at least $p+d$
exponents in residue zero and at least $d$ in a residue
$a\in\mathbf F_p^\times$.  Project its residue algebra onto
\[
 R'=\mathbf F_p[T]/T^{p+d}\ \times\
                    \mathbf F_p[T]/(T-a)^d.
\]
For $F=(F_0,F_a)$ in this quotient define
\begin{equation}\label{eq:sharp-functional}
 \begin{split}
 \lambda(F)={}&[Y^{d-1}]F_a(a+Y)
       -[T^{d-1}]F_0(T)-da[T^d]F_0(T)\\
       &\hspace{32mm}-a[T^{p+d-1}]F_0(T).
 \end{split}
\end{equation}
All indicated coefficients are well-defined in the truncated factors.
Set $C=p^2+d-1$.  Lucas's congruence gives
\begin{equation}\label{eq:sharp-functional-values}
 \begin{aligned}
 \lambda(T^k)&=0&&\text{if }s_p(k)<d,\\
 \lambda(T^k)&=0&&\text{if }s_p(k)=d\text{ and }k<C,\\
 \lambda(T^C)&=a\ne0.&
 \end{aligned}
\end{equation}
Here is a direct check.  The first term of
\eqref{eq:sharp-functional} on $T^k$ is
$\binom{k}{d-1}a^{k-d+1}$.  Since $d-1<p$, its binomial factor
vanishes unless the units digit of $k$ is at least $d-1$.
With digit sum less than $d$, the only exception is $k=d-1$,
cancelled by the second term.  With digit sum $d$ and $k<C$, the
only exceptions are $k=d$ and $k=p+d-1$, cancelled by the third and
fourth terms.  At $C$ the first term equals $a^{p^2}=a$, and the
other three terms vanish.  This verification also covers $d=1$.

Compose $\lambda$ with the residue projection and lift its values on
$1,T,\ldots,T^{s-1}$ to ordinary integers.  This defines an integral
functional $L$ on $A$, and hence rational coefficients $c_n$ by
Vandermonde inversion.  Its moments reduce to
\eqref{eq:sharp-functional-values}.  Choose, explicitly,
\[
 r=\frac1{p-1}+\frac1{2(p-1)C}.
\]
Then
$w_r(C)=(d+\tfrac12)/(p-1)<\min(1,(d+1)/(p-1))$.
All moments of smaller digit sum, and all preceding indices of digit
sum $d$, are divisible by $p$ and therefore have coefficient valuation
at least one.  Indices of greater digit sum have weight at least
$(d+1)/(p-1)$; later indices of the same digit sum have strictly
larger weight.  The $C$-th moment is a unit.  Thus $C$ is the unique
minimizing coefficient index, giving exactly $C$ zeros on this closed
disk, and therefore at least $C$ in $\mathcal D_p$.

For Theorem~\ref{thm:large-residue-class}, take
$d=\min(m-p,q)$ and choose a residue attaining $q$.  The quotient
above is available.  For Theorem~\ref{thm:sharp-staircase}, put
$d=\lfloor(t-p)/2\rfloor$ and use the support
\begin{equation}\label{eq:sharp-support}
 E_{p,d}=\{jp:0\le j<p+d\}\ \cup\
                   \{1+jp:0\le j<d\}.
\end{equation}
It has $p+2d$ exponents.  If this is $t-1$, add any new nonnegative
exponent.  The previously obtained root count survives a sufficiently
small nonzero coefficient on that exponent.

For completeness, both the full-support and simplicity requirements
can be enforced without relying on genericity.  First clear
denominators.  Perturb each permitted integer coefficient by a
multiple of a sufficiently large power $p^H$.  Each moment changes
by a number of valuation at least $H$, uniformly in its index, so
\eqref{eq:sharp-weight} shows that the strict minimizing index is
preserved.  Require all permitted coefficients to be nonzero.
After removing the smallest monomial, choose a prime $\ell\ne p$.
The Chinese remainder theorem additionally makes the leading
coefficient prime to $\ell$, all other coefficients divisible by
$\ell$, and the constant coefficient divisible by $\ell$ but not by
$\ell^2$.  The resulting polynomial is Eisenstein at $\ell$ and hence
irreducible over $\mathbf Q$.  In characteristic zero its roots are
simple.  When the original support contains a nonzero smallest
exponent, irreducibility is asserted after that monomial is removed;
the supports \eqref{eq:sharp-support} contain zero.  Finally, scaling
$X$ by $A$ proves the statement for an arbitrary center.

\subsection{Independent finite checks}

The script \texttt{code/verify\_sharp\_range.py} implements the residue
algebra directly by jets $\binom{k}{j}a^{k-j}$, using only arithmetic
in $\mathbf F_p$.  It scans powers by digit sum and then by index;
this is independent of the rational determinant implementation used
for the general algorithm.  It exhausts all residue-multiplicity
vectors at $p=5$ for $t=7,8,9$, checks larger-prime examples and
deterministically seeded profiles, and verifies
\eqref{eq:sharp-functional-values} separately.  For $p=5,t=9$ it also
constructs an integer Eisenstein witness, expands it directly at
$X=1$, and verifies that its Newton weights have last minimum at
$26$ at depth $53/208$.  These checks illustrate and test the argument;
the proofs above apply to every prime and every support in the stated
range.
